%% =====================================================================
%%  style-demo-body.tex —— 自动生成，请勿手改
%%  生成命令：cd preview/style-demo && python3 make-demo.py
%%  来源：Chapters/VolumeI/Chapter{2,7,13}.tex 原样片段（未删改原文）
%%  说明：readability.sty 的新宏是“插入式”标记；原版式编译时由
%%        style-demo.tex 中的降级宏丢弃/还原，正文文字不做任何删减。
%% =====================================================================

%% ===== 第二章：密集纯文字（§2.1 引言）=====
\setcounter{chapter}{1}
\chapter{基本物理}

\chapterintro{本章回答一个朴素的问题：\keyidea{物理学的“基本”到底指什么}？费曼用“看人下棋”的比喻说明：基本物理就是这盘宇宙大棋的规则。阅读时请抓住三件事：一、三种检验规则的思路；二、物理规律是逐层近似的；三、每次“综合”都带来新的问题。}
\readingpath{先读 \S 2.1 前四段建立直觉，再重点看“三种检验方法”，最后用页末小结回顾。}

\section{引言}

在本章中，我们将考察有关物理学的最\term{基本概念}——即我们在目前所知道的事物的本性。这里将不去涉及我们如何知道所有这些观念是正确的那个认识过程，你们在适当的时候会学习到这些具体的细节。

我们在科学上所关心的事物，具有无数形式和许多属性。举例来说，假如我们站在岸边眺望大海，将会看到：这里有海水、拍击的浪花、飞溅的泡沫以及汹涌的波浪，还有太阳、光线、蔚蓝的天空、白云以及空气的流动——风；在海边有沙粒，不同色纹和硬度的岩石；在海里浮游着生物，此生彼灭；最后，还有我们这些站在海岸边的观察者；甚至还有幸福和怀念。在自然界的其他场合，难道不也同样出现如此纷繁复杂的事物和影响吗？无论在哪里，到处都是这样错综复杂和变化无穷。\keyidea{好奇心驱使我们提出问题，把事物联系起来，而将它们的种种表现理解为：}或许是由较少量的基本事物和相互作用以无穷多的方式组合后所产生的结果。\aside{费曼在这里给本科生一个“世界观”：复杂现象可能由少量基本事物组合而成。}

例如，沙粒和岩石是两回事吗？就是说，沙粒只不过是大量的细小石块吗？月亮是不是一块巨大的岩石呢？如果我们了解岩石，是否就能了解沙粒和月亮呢？风是否与海洋中的水流相似，就是一种空气的流动？不同的运动有什么共同特征？不同的声音有什么相似之处？究竟有多少种颜色？等等，等等。我们就是试图这样地逐步分析所有的事情，把那些乍看起来似乎不相同的东西联系起来，希望有可能\uwave{减少不同类}事物的数目，从而能更好的理解它们。

几个世纪以前，人们想出了一种部分解答这类问题的方法，那就是：\uwave{观察}，\uwave{推理}和\uwave{实验}；这些内容构成了通常所说的\uwave{科学方法}。在这里，我们将只限于对那些有时称之为基本物理中的基本观点，或者由于应用科学方法而形成的基本概念作一描述。

现在我们要问：所谓“理解”某种事情指的是什么意思？可以作一想象：组成这个“世界”的运动物体的复杂排列似乎有点像是天神们所下的一盘伟大的象棋（这里指的是国际象棋——译者注），我们则是这盘棋的观众。我们不知道弈棋的规则，所有能做的事情就是\uwave{观看}这场棋赛。当然，假如我们观看了足够长的时间，总归能看出几条规则来，这些弈棋规则就是我们所说的基本物理。但是，即使我们知道了每条规则，仍然有可能不理解为什么下棋时要走某一步棋，这仅仅是因为情况太复杂了，而我们的智力却是有限的。如果你们会下棋，就一定知道，学会所有的规则是容易的，但是，要选择最好的一着棋，或者要弄懂别人为什么走这一着棋，往往就很困难了。在自然界里，也正是如此，而且只有更难一些。但是，至少我们能发现所有的规则。实际上我们今天还没找到一切规则（时而会出现一些像弈棋中“以车护王”那样的情况，使我们仍然感到无法理解）。除此之外，我们确实能用已知规则来解释的事情也是非常有限的，因为几乎所有的情况都是极其复杂的，我们不能领会这盘棋中应用这些规则的走法，更无法预言下一步将要怎样。所以，我们必须使自己只限于弈棋规则这个比较基本的问题。如果我们知道了规则，就认为“理解”了世界。

\runin{三种检验方法}如果我们不能很好的分析这盘象棋游戏，那么又怎样来辨别我们“猜测”出的规则实际上是否正确呢？大致地讲，可以有三种办法。第一，可能有这种情况：大自然安排的，或者说我们将大自然安排的十分简单，只有少数几个组成部分，从而使我们能够正确地预测将要发生的事。在这种情况下，就能检验我们的规则是怎样起作用的。（在棋盘角落里可能只有少数几个棋子在移动，所以我们能够正确地解决。）

第二种检验规则的好办法是，利用那些由已知规则推导出来的一些较一般性的法则来检验已知规则本身。比如，象在棋盘中移动的规则是只许走对角线，因而我们可以推断，无论象走了多少步，它总是出现在红方块里。这样，即使不能领会细节，我们也总能检验有关象的走法的概念，只要弄清楚它是否一直在红方块里。当然，在相当长的时间里，它都将如此，直到突然发现它出现在黑方块里。（显然，这时发生的情况是这个象被俘获了，另一个卒走过来成为皇后，红方块里的象就变成黑方块里的象。）这也就是物理学中出现的情况，即使我们不能领会其中的细节，但是在相当长的时期内我们仍有在各方面都很好地起作用的规则；但是在某个时候，我们又会发现新的规则。从基本物理的观点来看，最有趣的现象当然是在那些\uwave{新的}场合——那些已知规则行不通的场合中所出现的现象，而不是在原有规则行得通的地方发生的现象！这是我们发现新规则的一条途径。

第三个鉴别我们的观念是否正确的方法比较粗糙，但或许是所有方法中最为有效的。这就是用粗略的近似方法来加以辨别。我们可能说不出为什么阿莱克因（Alekhine）\footnote{世界著名弈棋手，系国际象棋大师。曾多次获得国际象棋世界冠军。——译者注}要走这步棋，但是我们或许能大致认为它或多或少地在调集一些棋子到王的周围来保护它。因为这是在这种情况下明摆着的事。同样，根据我们对这盘棋的理解，即使不能看出每一步棋的作用，也常常能对自然界多少有所理解。

\begin{summarybox}[本节要点]
  \begin{itemize}
    \item 基本物理研究的是“规则”：用少量基本事物与相互作用解释纷繁现象。
    \item 检验规则有三条路：预测简单情形；用推论检验规则本身；用粗糙近似作判断。
    \item 已知规则只是近似，新场合（尤其规则失效处）才是发现新规则的地方。
  \end{itemize}
\end{summarybox}

\clearpage

%% ===== 第七章：图表密集（§7.1--§7.3）=====
\setcounter{chapter}{6}
\chapter{万有引力理论}

\chapterintro{引力是第一个被“统一”的力：一条 $F=Gmm'/r^2$ 就能同时解释落体、月球、行星与潮汐。本章跟着费曼走一遍“观测数据 $\to$ 定律 $\to$ 动力学解释”的推理链。}

\section{行星运动}
\readingpath{重点看开普勒三定律如何从第谷的观测数据中被“读”出来，以及牛顿如何用动力学解释它们。}

在这一章中，我们将要讨论对人类智慧影响至为深远的概括之一的引力定律。当我们现在赞颂人类智慧的时候，应当先停下来向\uwave{大自然}表示敬畏之意，因为她能如此完整而普遍地遵循引力定律这样一个出奇地简单的原理。那么什么是\term{引力定律}呢？它指出，宇宙中每一个物体都以一定的力吸引着每一个其他物体，而对任何两个物体来说，这一力正比于每一个物体的质量，而反比于它们之间距离的平方。这个称述数学上可以用下列式子来表示：
\begin{equation*}
F=G\,\frac{mm'}{r^2}.
\end{equation*}
如果对此再加上一个事实，即一个物体在力的作用下会沿着力的方向得到加速，而加速的快慢与物体的质量成反比；那么我们就已说出了所需要的一切，于是一个天资卓越的数学家就能推导出这两个原理的所有结论。然而由于你们还没有被认为天资如此卓越，所以我们要更详细地来讨论一下这些结论，而不是只给你们留下简单的原理。我们将简短地叙述一下发现引力定律的故事，讨论它的某些结果，它在历史上的作用，这样一条定律所遗留下来的神秘之处，以及爱因斯坦对这条定律所作的若干改进；我们还将讨论这条定律与物理学中其他定律的关系。所有这些不可能在一章中都讲到，所以有些论题将在适当的时候放到稍后的几章中去讨论。

故事要从古人对行星在恒星中间运动的观察，并且最终作出了它们在围绕太阳运行的推论开始，这是后来为哥白尼所重新发现的一个事实。行星究竟\uwave{怎样}围绕太阳运行，并且究竟用什么样的\uwave{运动}绕之运行，要发现这些，就要稍微多作一点工作。15世纪初叶，在行星到底是不是围绕太阳运行这个问题上曾有过激烈的争论。第谷•布拉赫(Tycho Brahe)有一个想法，它与古人提出的任何观点都不相同，他认为：如果能足够精确地测得行星在天空中实际的位置，那么这些有关行星运动本性的争论就会得到最好的解决。如果测量能精确地显示出行星在如何运动，那么或许有可能去建立这种或那种观点。这是一个非同小可的想法：\keyidea{如果要想发现什么东西，那么去细致地做一些实验要比展开冗长的哲学争辩好得多。}在这个想法的指引下，第谷•布拉赫在哥本哈根附近的希恩(Hven)岛上他的天文台里，花了多年时间来研究行星的位置。他编制了一种篇幅庞大的星表；在第谷死后，数学家开普勒对这些星表进行了研究，从这些数据中，开普勒发现了涉及行星运动的一些非常优美、卓越而又简单的定律。\aside{第谷的观测精度是开普勒定律的前提：没有好数据，就没有好定律。}



\section{开普勒定律}

开普勒首先发现，每个行星沿一条称为\uwave{椭圆}的曲线绕太阳运行，而太阳处在椭圆的一个焦点上。椭圆不仅仅只是一个卵形的东西，而是一条非常独特和精确的曲线，这条曲线可以用两只平头钉（在每个焦点上各钉一只），一束线和铅笔把它画出来；或者用数学术语来说，椭圆是（平面上）到两个定点（焦点）的距离之和是一个常数的轨迹。或者，如果你愿意的话，就把它说成是一个压扁了的圆吧（图7.1）。

\begin{wrapfigure}{r}{0.3\textwidth}
    \centering
    \includegraphics[width=0.25\textwidth]{Chapter7/椭圆}
    \caption{椭圆\figsource{《费曼物理学讲义》第一卷 §7-2}}
    \label{figure:椭圆}
\end{wrapfigure}
其次开普勒发现，行星并不以均匀速率绕太阳转动，而是当它们接近太阳时跑得较快，远离太阳时则跑得较慢，确切地说便是这样：设在任意相继的两个时间，比如说相隔为一周的时间内观察一个行星，并且对每个观察位置向行星画一条矢径\footnote{矢径是从太阳到行星轨道上任何一点的连线。}。那么行星在一周中所经过的轨道上一段弧线和两条矢径一起围成一定的平面面积，如图7.2所示的那个阴影面积。如果在离太阳较远的那部分轨道上（此时行星运动得较慢），也作时间相隔一周的与前类似的两次观察，那么这时围成的面积与前一情况下的面积完全相等。因此，按照开普勒第二定律，每个行星的轨道速率都使矢径在相等的时间内“扫过”相等的面积。

\begin{wrapfigure}{l}{0.3\textwidth}
    \centering
    \includegraphics[width=0.25\textwidth]{Chapter7/开普勒的面积定律}
    \caption{开普勒的面积定律\figsource{《费曼物理学讲义》第一卷 §7-2}}
    \label{figure:开普勒的面积定律}
\end{wrapfigure}
\runin{开普勒第三定律}开普勒第三定律发现得较晚，这条定律与前两条不同，各属于不同的范畴，因为它不是只涉及单独的行星，而涉及一个行星与其他行星之间的关系。这条定律表明：如果把任何两个行星的轨道周期和轨道大小进行比较，则周期与轨道大小的3/2次方成正比。这里所说的周期是行星在其轨道上完全绕一圈所需的时间间隔，而所谓轨道的大小使用椭圆轨道最大直径（术语叫“长轴”）的长度来量度的。更简单一些，如果行星绕圆周运动（实际上，它们也近于这样做），那么绕圆周走一圈所需的时间将正比于直径（或半径）的3/2次方。这样，开普勒三条定律便是：
\begin{spacing}{0.3}
\begin{enumerate}
\item[\uppercase\expandafter{\romannumeral1}.] 每个行星都沿椭圆轨道绕太阳运行，太阳位于各椭圆的其中一个焦点上。
\item[\uppercase\expandafter{\romannumeral2}.] 从太阳画到行星的矢径，在相等时间间隔内扫过相等的面积。
\item[\uppercase\expandafter{\romannumeral3}.] 任何两个行星的周期平方正比于它们各自轨道半长轴的立方：$T^2\propto a^3$（或 $T\propto a^{3/2}$）。
\end{enumerate}
\end{spacing}

\section{动力学的发展}

当开普勒发现这些定律的时候，伽利略正在研究有关运动的定律。当时的问题在于什么东西使行星绕转运动（那时有一种理论这样说，行星之所以运行是因为在它们背后有看不到的天使在扑动它的飞翼，推动它们前进。你们将会看到，这个理论现在被修改了一下！这就是说，为了保持行星的绕转运动，看不见的天使必须朝与运动方向不同的方向飞行，并且它们也没有飞翼。除此之外，它倒多少有点像先前的理论！）在有关运动方面，伽利略发现了一个非常值得注意的事实，这个事实对于理解开普勒定律是必不可少的。这就是\uwave{惯性}原理——如果有某个物体在运动，但没有和其他东西相碰撞，也完全不受干扰，那么它将沿一直线以均匀速度永远运动下去。（\uwave{为什么}它能保持直线运动？我们不知道，但是事情就是如此。）

牛顿使这个观念更为明确，他说：“改变物体运动的唯一方法是要对之用\uwave{力}。”如果物体的速率变大，这必定有一个力施加在\uwave{运动方向上}。另一方面，如果物体的运动改变到另一个新的方向，那么它必定受到一个\uwave{斜向}的力作用。这样牛顿添进了如下一个概念：要改变一个物体运动的速率\uwave{或方向}，就需要有力才行。例如：把一块石子系在绳上，使它旋转而作圆周运动，那么就需要有一个力以保持它在圆周上运行。这时我们必须把绳子\uwave{拉}住。事实上，这个定律说的是，力所产生的加速度反比于物体的质量；或者说，力正比于质量乘加速度。物体的质量越大，使它产生某一给定加速度所需的力就越大。（质量可以这样来测量，使其他石子系于同一根绳的末端，使它们以同样的速率绕同样的圆周转动，用这种方法可以知道它们所需的力的大小，质量较大的物体，所需的力较大。）从这些考虑中得出的一个卓越的观念就是：要保持行星在它的轨道上运行，根本不需要有一个\uwave{切向的力}（天使并不一定要沿切线方向飞行），因为行星总会沿所要求的方向运动，如果根本没有什么东西去干扰它，那么行星就将沿\uwave{直线}运行下去。但实际的运动却偏离了不存在力作用时物体所应沿之运动的那条直线，这种偏差差不多与运动相垂直，而不沿运动的方向。换句话说，由惯性原理得知，控制行星绕太阳运动所需的力不是一个\uwave{绕}太阳而是\uwave{指向}太阳的力。（如果有一个力指向太阳，那么当然太阳也许就是那天使了！）

\begin{summarybox}[本节要点]
  \begin{itemize}
    \item 开普勒三定律把“行星怎么动”总结成三条精确而简单的规律。
    \item 牛顿用“力指向太阳 + 反平方”解释了面积定律与第三定律。
    \item 引力定律的普适性来自“天上的力”与“地上的力”被认定为同一种力。
  \end{itemize}
\end{summarybox}

\clearpage

%% ===== 第十三章：公式密集（§13.1 落体的能量）=====
\setcounter{chapter}{12}
\chapter{功与势能（上）}

\chapterintro{本章用最朴素的例子把“动能 + 势能 = 恒量”推出来：先看自由落体，再推广到沿曲线运动、再到三维。公式不多，关键是理解每一步“求导—代牛顿定律—识别全导数”的结构。}

\section{落体的能量}
\readingpath{先看 $T+U=$ 恒量的结论，再跟上两条求导推导；看不懂推导时，先记住“动能变化率 = 力 × 速度”这个核心。}

在第4章中我们讨论过能量守恒。那里，我们没有引用牛顿定律，但是，按照牛顿定律，能量实际上是守恒的。我们来看一看能量守恒如何与牛顿定律相符合是很有意义的。为了清楚起见，我们从最简单的实例开始讨论，然后逐步推广到比较难的例子。

能量守恒最简单的例子是一个垂直下落的物体，它只在垂直方向上运动。一个仅仅在重力作用下改变高度的物体，由于下落运动而具有\term{动能}$T$（\enterm{kinetic energy}，或K.E.），并且还具有势能$mgh$，简写成$U$（或P.E.）。\keyidea{这两种能量的总和是恒量}
\begin{equation*}
    \mathop{\frac{1}{2}mv^2}\limits_{K.E.}+\mathop{mgh}\limits_{P.E.}=\text{恒量}
\end{equation*}
或
\begin{equation}
    \label{Eq:I:13:1}
    T+U=\text{恒量}
\end{equation}
现在我们要证明这一表述是正确的。我们说“证明这一表述是正确的”是什么意思呢？从牛顿第二定律我们很容易说明物体如何运动，而且容易求出速度如何随时间而变化——速度的增加与时间成正比，高度随时间的平方而改变。所以，假若我们以物体静止的那个位置作为零点来测量高度，那么，高度等于速度的平方乘以一些常数就不足为奇了。不管怎样，我们还是来稍微仔细地看一下这个问题。

我们将动能对时间求微商，然后应用牛顿定律直接从第二定律求出动能是如何变化的。我们把$mv^2/2$对时间求微商，由于假设$m$是常数，故得
\begin{equation}
    \label{Eq:I:13:2}
    \frac{dT}{dt}=\frac{d(\frac{1}{2}mv^2)}{dt}=\frac{1}{2}m \cdot 2v \frac{dv}{dt}=mv\frac{dv}{dt}
\end{equation}
但由牛顿第二定律，$m(dv/dt)=F$，所以
\begin{equation}
    \label{Eq:I:13:3}
    \frac{dT}{dt}=Fv
\end{equation}
一般地讲，结果应该是$\boldsymbol{F} \cdot \boldsymbol{v}$，但在一维情况下我们只要写成力乘以速度就够了。

在我们的简单例子中，力是常数，等于$-mg$，即一个垂直向下的力（负号的意思是指作用向下），而速度当然是垂直位置（或高度）对时间的变化率。这样，动能的变化率就是$-mg(dh/dt)$，非常奇怪，这个量是另外的物理量的变化率的负数，它是$mgh$的时间变化率的负数！因此，随着时间的推移，动能的变化在数量上等于$mgh$的变化，而符号相反；所以，这两个量的总和保持不变，证明完毕。

从牛顿第二定律我们已经证明，在恒力的情况下如果将势能$mgh$与动能$mv^2/2$加在一
\begin{wrapfigure}{r}{0.4\textwidth}
    \centering
    \includegraphics[width=0.35\textwidth]{Chapter13/物体在重力作用下在一无摩擦的曲线上的运动}
    \caption{物体在重力作用下在一无摩擦的曲线上的运动\figsource{《费曼物理学讲义》第一卷 §13-1}}
    \label{figure:物体在重力作用下在一无摩擦的曲线上的运动}
\end{wrapfigure}
起，则能量是守恒的。现在我们进一步看一下是否能将它推广，以加深我们的理解。能量守恒定律是否只对自由落体适用，还是可以适用于更一般的情况？根据对能量守恒的讨论，我们预计它对一个物体在重力作用下沿一曲线无摩擦地从一点运动到另一点的情形也是成立的（图\ref{figure:物体在重力作用下在一无摩擦的曲线上的运动}）。

如果物体从原来的高度$H$到达某一高度$h$，即使速度不再是沿垂直方向，上述方程仍然应当成立。我们要搞清楚为什么这条定律仍然正确。下面我们用同样的分析方法求出动能的时间变化率。诚然，动能的时间变化率仍是$mv(dv/dt)$，而$m(dv/dt)$则是动量大小的变化率，也就是运动方向上的力——切向力$F_t$。于是有
\begin{equation*}
    \frac{dT}{dt}=mv\frac{dv}{dt}=F_tv
\end{equation*}
这里的速率是沿着曲线的距离对时间的变化率$ds/dt$，但切向力$F_t$不是$mg$，而是随垂直的距离$dh$与沿路径的距离$ds$的比率而减弱，换句话说
\begin{equation*}
    F_t=-mg\sin \theta=-mg\frac{dh}{ds}
\end{equation*}
所以
\begin{equation*}
    F_t\frac{ds}{dt}=-mg(\frac{dh}{ds})(\frac{ds}{dt})=-mg\frac{dh}{dt}
\end{equation*}
因为$ds$被消掉了，于是我们就得到$-mg(dh/dt)$，与前面所证明的一样，它等于$-mgh$的时间变化率。

为了确切地理解能量守恒定律一般在力学中是怎样起作用的，我们来讨论几个有助于分析这个问题的概念。

首先我们讨论三维情况下一般的动能变化率。在三维空间的动能是
\begin{equation*}
    T=\frac{1}{2}m(v_x^2+v_y^2+v_z^2)
\end{equation*}

\begin{summarybox}[本节要点]
  \begin{itemize}
    \item 能量守恒最简单的形式：$T+U=\text{恒量}$。
    \item 证明路径：对 $T$ 求时间导数，代入 $F=ma$，认出结果恰是 $-U$ 的导数。
    \item 该结论可推广到沿任意曲线运动，只要力是保守力。
  \end{itemize}
\end{summarybox}
